%! AP Statistics — Unit 2, Lesson 2.1 — Group Activity KEY
\documentclass[10pt]{article}
\usepackage{apstats-article}
\usepackage{apstats-key}

%=============================================================================
\begin{document}

\pageheader{Unit 2, Lesson 2.1}{Group Activity --- Answer Key}
\begin{tabularx}{\linewidth}{X X X}
Name:\;\ans{Key} & Date:\; & Period:\;
\end{tabularx}

\vspace{0.1in}

\begin{tcolorbox}[colback=sky, colframe=navy, boxrule=0.8pt, arc=2mm,
    left=3mm, right=3mm, top=2mm, bottom=2mm]
\small
For each study below: (1) identify the two variables and their types,
(2) decide whether the apparent association is likely \textbf{real} or possibly \textbf{random},
(3) name one possible \textbf{confounding variable}, and (4) state whether the study
could establish \textbf{causation}.
\end{tcolorbox}

\vspace{0.14in}

\begin{scenariobox}[Scenario 1]{navylight}

\begin{headlinebox}{sky}
\small\textit{``A study of 1{,}500 adults found that those who own pets report higher
levels of happiness than those who do not own pets.''}
\end{headlinebox}

\vspace{8pt}
\renewcommand{\arraystretch}{1.8}
\begin{tabularx}{\linewidth}{@{} l X @{}}
\textbf{Two variables:}          & \ans{Pet ownership; happiness level} \\
\textbf{Variable types:}         & \ans{Pet ownership: categorical (yes/no); happiness: quantitative (or categorical if scored)} \\
\textbf{Real or possibly random?}  & \ans{Likely real --- $n = 1{,}500$ is large enough that a random pattern is unlikely} \\
\textbf{One confounding variable:} & \ans{Accept any reasonable answer: e.g., income (wealthier people can afford pets and also tend to report higher happiness); living alone vs.\ with others; physical activity} \\
\textbf{Establishes causation?}  & \ans{No --- observational study; we cannot randomly assign pet ownership} \\
\end{tabularx}

\end{scenariobox}

\vspace{0.16in}

\begin{scenariobox}[Scenario 2]{greenacc}

\begin{headlinebox}{greenbg}
\small\textit{``Researchers found that states with higher ice cream sales also have
higher drowning rates.''}
\end{headlinebox}

\vspace{8pt}
\renewcommand{\arraystretch}{1.8}
\begin{tabularx}{\linewidth}{@{} l X @{}}
\textbf{Two variables:}          & \ans{Ice cream sales; drowning rate} \\
\textbf{Variable types:}         & \ans{Both quantitative} \\
\textbf{Real or possibly random?}  & \ans{The association is likely real (consistent across many states) but non-causal} \\
\textbf{One confounding variable:} & \ans{Temperature / season: hot weather causes both more ice cream sales and more swimming, leading to more drownings} \\
\textbf{Establishes causation?}  & \ans{No --- ice cream does not cause drowning; temperature (or summer season) is the confounding variable} \\
\end{tabularx}

\vspace{8pt}
\ans{The association is real because temperature causes both. Hot weather makes people buy more ice cream \emph{and} go swimming more (leading to more drownings). Ice cream is not the cause.}

\end{scenariobox}

\vspace{0.16in}

\begin{scenariobox}[Scenario 3 \normalfont\small\color{white!80!black}\quad stretch]{redacc}

\begin{headlinebox}{redbg}
\small\textit{``A study of 12 high school students found that the 4 students who slept
the most also had the highest grades. The 4 students who slept the least had the lowest grades.''}
\end{headlinebox}

\vspace{8pt}
\renewcommand{\arraystretch}{1.8}
\begin{tabularx}{\linewidth}{@{} l X @{}}
\textbf{Two variables:}          & \ans{Hours of sleep; GPA (or grade level)} \\
\textbf{Variable types:}         & \ans{Both quantitative} \\
\textbf{Real or possibly random?}  & \ans{Possibly random --- with only 12 students, this pattern could easily occur by chance} \\
\textbf{One confounding variable:} & \ans{Time management / academic motivation: students who are more organized may both sleep better and study more effectively} \\
\textbf{Establishes causation?}  & \ans{No --- even if real, this is observational; other variables could explain both sleep and grades} \\
\end{tabularx}

\vspace{8pt}
\ans{With only 12 students, a pattern this strong could easily be due to random chance.
If we sampled a different 12 students, the pattern might look completely different.
We need a much larger sample (and statistical tools from Units 5--9) before we can
be confident the association is real.}

\end{scenariobox}

\vspace{0.16in}

\begin{tcolorbox}[colback=goldbg, colframe=goldacc, boxrule=0.8pt, arc=2mm,
    left=3mm, right=3mm, top=2mm, bottom=2mm]
\small\textbf{Reflection:}\\[4pt]
\ans{Scenario 2 (ice cream and drowning) best illustrates association $\neq$ causation.
The association is large and consistent, yet no one would claim ice cream causes drowning.
The confounding variable (temperature/season) explains both, showing that even a strong
and real association can be entirely non-causal.}
\end{tcolorbox}

\vspace{0.12in}
\begin{teachernote}
\textbf{Scenario 1:} Multiple valid confounders exist (income, social connection, lifestyle).
Accept any that are related to \emph{both} pet ownership and happiness.

\textbf{Scenario 2:} This is the canonical ``spurious correlation'' example.
Push students to name the mechanism: hot weather $\to$ more ice cream, hot weather $\to$
more swimming $\to$ more drowning. The key insight is that both variables are effects
of a common cause.

\textbf{Scenario 3:} The stretch purpose is to highlight sample size.
Students who say ``the association could be random'' are correct --- that's the main point.
If students say ``it proves sleep causes better grades'' they need to be redirected on two fronts:
(1) small sample, and (2) observational study.
\end{teachernote}

\end{document}
